\documentclass[11pt]{article}
\usepackage[T1]{fontenc}
\usepackage[utf8]{inputenc}
\usepackage{lmodern}
\usepackage[ngerman]{babel}
\usepackage{amsmath,amssymb,mathtools,array,booktabs}
\usepackage{geometry}
\usepackage{microtype}
\geometry{a4paper,margin=1.45cm}
\setlength{\parindent}{1.25em}
\setlength{\parskip}{0pt}
\allowdisplaybreaks
\setlength{\emergencystretch}{10em}
\sloppy
\hbadness=10000
\vbadness=10000
\hfuzz=2pt
\vfuzz=10pt
\raggedbottom
\providecommand{\frR}{\mathfrak R}
\providecommand{\frT}{\mathfrak T}
\providecommand{\frM}{\mathfrak M}
\providecommand{\frN}{\mathfrak N}
\providecommand{\frO}{\mathfrak O}
\providecommand{\frS}{\mathfrak S}
\providecommand{\frB}{\mathfrak B}
\providecommand{\frC}{\mathfrak C}
\providecommand{\frA}{\mathfrak A}
\providecommand{\frakp}{\mathfrak p}
\providecommand{\ideal}[1]{\mathfrak{#1}}
\providecommand{\qmod}[1]{\pmod{#1}}
\begin{document}

\section*{30. Abstrakter Aufbau der Idealtheorie in algebraischen Zahl- und Funktionenkörpern}
\noindent\emph{Math. Ann. 96 (1927), S. 26--61.}

Im folgenden wird eine abstrakte Charakterisierung all derjenigen Ringe gegeben,
deren Idealtheorie übereinstimmt mit der Idealtheorie aller ganzen Größen des
algebraischen Zahlkörpers — deren Ideale sich also eindeutig als Potenzprodukte
von Primidealen darstellen lassen. Zu diesen Ringen gehören als bekannteste
Beispiele noch der Ring aller ganzen Größen eines algebraischen Funktionenkörpers
von einer Unbestimmten — allgemeiner von mehr Unbestimmten, sobald man sich
mit Kronecker auf Ideale höchster Dimension beschränkt, also den Übergang zu
einem gewissen Quotientenring, dem Funktionalbereich, macht.

Die Axiome, die dieser Idealtheorie vollständig äquivalent sind, sind für den
zugrunde gelegten kommutativen Ring die folgenden:\footnote{Für die Definition
der Grundbegriffe der Idealtheorie vgl. etwa E. Noether, \emph{Idealtheorie in
Ringbereichen}, Math. Ann. 83 (1921), S. 24--66 (zitiert: \emph{Idealtheorie}).
Übrigens werden die wesentlichsten Begriffe auch in der vorliegenden Arbeit,
besonders in §§1, 4 und 5, kurz formuliert. Der Vollständigkeit halber ist dabei
manches aufgenommen, was für die unmittelbaren Zwecke dieser Arbeit nicht
erforderlich ist.}

\begin{description}
\item[I. Teilerkettensatz.] Jede Kette von Idealen, bei der jedes Ideal ein
echter Teiler des vorangehenden ist, bricht im Endlichen ab — m.\,a.\,W.:
zu jeder Teilerkette von Idealen gibt es einen Index, von dem an alle Ideale
gleich werden.
\item[II. Vielfachen-Kettensatz modulo jedem vom Nullideal verschiedenen Ideal.]
Jede Kette von Idealen — die sämtlich Teiler eines festen, vom Nullideal
verschiedenen Ideals sind —, bei der jedes Ideal ein echtes Vielfaches des
vorangehenden ist, bricht im Endlichen ab.
\item[III.] Existenz des Einheitselementes der Multiplikation.
\item[IV.] Ring ohne Nullteiler.
\item[V. Ganze Abgeschlossenheit im Quotientenkörper.] Jedes Element des
Quotientenkörpers, das ganz in bezug auf den Ring ist, gehört dem Ring an.
\end{description}

Aus diesen Axiomen entwickle ich schrittweise eine immer stärker eingeschränkte
Idealtheorie bis hin zu der gesuchten; und zeige umgekehrt, daß hier auch alle
Axiome tatsächlich erfüllt sind.

Aus dem Teilerkettensatz I folgt, wie ich (\emph{Idealtheorie} §4) gezeigt habe,
die Darstellung eines Ideals als kleinstes gemeinsames Vielfaches von endlich
vielen, zu verschiedenen Primidealen gehörigen Primäridealen. Ich wiederhole
hier kurz diesen Beweis (§6), weil er sich unter Verwendung des Begriffs des
Idealquotienten vereinfachen läßt, und weil ich ein Versehen berichtigen will:
der ursprüngliche Beweis benutzt nämlich an einer Stelle die nicht vorausgesetzte
Existenz des Einheitselementes der Multiplikation.\footnote{Auf dieses Versehen
machte mich P. Urysohn aufmerksam; meine Abänderung des Beweises war aber
etwas umständlicher als die hier mitgeteilte, die von E. Artin stammt. Dieser
hatte schon früher für sich die Lücke ausgefüllt und sich auch unabhängig von mir
die Vereinfachung mit dem Idealquotienten überlegt. — Eine weitere Vereinfachung,
die die Einführung der „reduzierten Darstellung“ vermeidet, entnehme ich einer
verwandten Fragestellung bei W. Krull: \emph{Algebraische Theorie der Ringe III},
Math. Ann. 92 (1924), S. 181--213, Satz I.}

Die Voraussetzung des Vielfachen-Kettensatzes bewirkt (§7), daß im
Restklassenring nach jedem vom Nullideal verschiedenen Ideal ein Primideal keinen
echten Teiler besitzen kann, mit Ausnahme des alle Elemente umfassenden
Einheitsideales. Damit aber werden die Primärkomponenten dieser Ideale eindeutig
bestimmt, mit Ausnahme der zum Einheitsideal gehörigen. In etwas weniger
scharfer Fassung findet sich dies Resultat schon bei Masazo Sono;\footnote{Masazo
Sono, \emph{On the Reduction of Ideals}, Memoirs of the College of Science,
Kyoto University, Series A, 7 (1924), S. 191--204.} seine Voraussetzung der
Existenz einer Kompositionsreihe im Restklassenring ist dem „Doppelkettensatz“
in diesem Ring gleichwertig, wie ich zum Schluß (§10) kurz zeige. Dabei verstehe
ich unter Doppelkettensatz die Gültigkeit von Teilerketten- und Vielfachenkettensatz
ohne Beschränkung auf vom Nullideal verschiedene Ideale.

Ist noch Axiom III — Existenz des Einheitselementes — erfüllt, so tritt das
Einheitsideal nicht als zugehöriges Primideal auf; die Primärkomponenten sind also
eindeutig bestimmt; sie werden paarweise teilerfremd und ihr kleinstes gemeinsames
Vielfaches gleich dem Produkt. Die Axiome I, II, III sind für alle endlichen
Ordnungen eines algebraischen Zahlkörpers erfüllt; hier hat schon Dedekind
(Dirichlet--Dedekind, \emph{Zahlentheorie}, 3. Aufl., 11. Suppl., §172) ohne
vollständig ausgeführten Beweis die eindeutige Darstellung eines Ideals als Produkt
von paarweise teilerfremden einartigen Idealen (Primärkomponenten) ausgesprochen.

Aus Voraussetzung IV — Nichtauftreten von Nullteilern — folgt, daß die
verschiedenen Potenzen eines von Null- und Einheitsideal verschiedenen Ideals alle
verschieden sind, und daß der Quotientenkörper existiert. Ist schließlich noch
Axiom V der ganzen Abgeschlossenheit im Quotientenkörper erfüllt, so werden die
Primärideale — wegen IV eindeutig bestimmte — Potenzen von Primidealen,
womit der übliche Zerlegungssatz gewonnen ist (§8). Dieser letztere Nachweis beruht
auf einer, im Fall des Zahlkörpers schon von Dedekind (3. Aufl., §172) gezogenen
Folgerung aus der ganzen Abgeschlossenheit: Man kann ein echt gebrochenes Element
so als Quotient ganzer darstellen, daß auch das Quadrat des Zählers nicht durch
den Nenner teilbar wird. An Stelle dieser Folgerung tritt in den späteren
Darstellungen — auch als Folge der ganzen Abgeschlossenheit — der viel weniger
durchsichtige verallgemeinerte Gaußsche Satz oder entsprechende, etwas mehr
aussagende Modulsätze; alles Sätze, die zur Anwendung Basiselemente voraussetzen,
während hier mit den Idealen selbst gearbeitet wird.

In der Umkehrung (§9) ergeben sich die von Axiom V verschiedenen Axiome fast
direkt; letzteres aus dem Nachweis, daß ein echt gebrochenes Element sich stets so
darstellen läßt, daß keine Potenz des Zählers durch den Nenner teilbar wird, was
also noch eine Verschärfung der ursprünglich aus V gezogenen Folgerung bedeutet.
Dieser Nachweis gelingt aber erst, indem aus der Idealdarstellung die üblichen
Schlüsse gezogen werden, Hauptidealdarstellung modulo einem festen Ideal und
Theorie der gebrochenen Ideale, also der Idealquotienten im Körper. Auch die
Tatsache, daß aus der Darstellung durch Primidealpotenzen die ganze
Abgeschlossenheit folgt, war schon Dedekind bekannt, wie eine Bemerkung
(3. Aufl., §172, S. 522 unten) zeigt.

Die Axiome I bis V ergeben, daß die vier im allgemeinen getrennten Zerlegungen
in kleinste paarweise teilerfremde Ideale, gegenseitig prime Ideale, größte
Primärkomponenten und irreduzible Ideale zusammenfallen; und umgekehrt ergibt
sich aus diesem Zusammenfallen und den Axiomen I--IV die ganze Abgeschlossenheit.
Für Ringe mit Nullteilern folgt aus dem Erfülltsein der übrigen Axiome — wobei
V als ganze Abgeschlossenheit im Quotientenring zu definieren ist — nicht notwendig
das Zusammenfallen der Zerlegungen, wie das Beispiel des Restklassenringes nach
gewissen Idealen einer endlichen Ordnung zeigt (§9).

In §1 skizziere ich die Theorie der in bezug auf einen Ring ganzen Größen, bis hin
zu der Dedekindschen Folgerung aus der ganzen Abgeschlossenheit. In §2 zeige ich,
wie die Kettensätze sich vom Ring übertragen auf Moduln aus Linearformen, mit
diesem Ring als Koeffizienten- und Multiplikatorenbereich.\footnote{Diese Fassung
der Modulsätze rührt von Prüfer her (Neue Begründung der algebraischen
Zahlentheorie, §2, Math. Ann. 94 (1924), S. 198--243) und ist durchsichtiger
als die übliche Übertragung der Basissätze.} Hieraus ziehe ich (§3) die Folgerung,
daß beim Übergang zu den ganzen Größen eines algebraischen Erweiterungskörpers
(erster Art) die Axiome I bis IV für jede Ordnung erhalten bleiben, während für
die aus allen ganzen Größen bestehende Ordnung auch V gilt. Diese — für den
algebraischen Zahlkörper natürlich bekannte — Tatsache erlaubt die Unterordnung
der am Anfang erwähnten Funktionenkörper; insbesondere ist durch den einfachen
Nachweis, daß der aus den ganzzahligen Polynomen mehrerer Unbestimmter abgeleitete
Funktionalbereich den Axiomen genügt, auch die Kroneckersche Idealtheorie voll
eingeordnet.\footnote{Die Verschiedenheit tritt also erst bei den Diskriminantensätzen
auf, dadurch bedingt, daß der Restklassenring des Funktionalbereiches nach einer
Primzahl ein unvollkommener Körper wird, und somit Erweiterungen zweiter Art
auftreten können.}

Die dann entwickelte Idealtheorie setzt diese nur der Einordnung dienenden §§2
und 3 nicht voraus. In §4 und §5 werden die von den Axiomen I bis V unabhängigen
Grundlagen gegeben; dadurch tritt schärfer als in der ursprünglichen Begründung
hervor, welche Teile der Theorie von Endlichkeitsvoraussetzungen unabhängig sind.

\subsection*{§1. Theorie der ganzen Größen}

Zugrunde gelegt sei ein kommutativer Ring\footnote{Ein Ring ist bekanntlich
definiert, wenn in einer Menge eine Gleichheitsrelation gegeben ist, und außerdem
zwei Verknüpfungen, Addition und Multiplikation, die den üblichen Gesetzen genügen:
die Addition ist assoziativ, kommutativ und eindeutig umkehrbar, die Multiplikation
ist assoziativ — kommutativ, wenn es sich um einen kommutativen Ring handelt —
und gegenüber der Addition distributiv. Ist dabei die zugrunde gelegte
Gleichheitsdefinition nicht die mengentheoretische Identität, so muß — damit in
jedem Untersystem dieselbe Gleichheitsdefinition erhalten bleibt — ein solches
Untersystem (Unterring, Modul, Ideal) neben irgendeinem Element auch alle ihm
gleichen enthalten. Ebenso muß bei jeder Erweiterung vorausgesetzt werden, daß
die neue Gleichheitsdefinition die alte umfaßt. Alle Begriffe wie „endlich viele
Elemente“, „eindeutige Zuordnung“, \dots{} sind im Sinne der Gleichheitsdefinition
gemeint; es gibt „ein und nur ein Element“ heißt also beispielsweise, es gibt bis
auf gleiche nur ein Element. Übrigens kann man von der Gleichheit immer zu einer
mengentheoretischen Identität übergehen, indem man gleiche Elemente in eine Klasse
zusammenfaßt und diese Klassen als neue Elemente einführt. Die hier gegebenen
Bemerkungen zur Gleichheitsdefinition stammen von R. Hölzer.} $T$ ohne Nullteiler,
mit Einheitselement $e$ der Multiplikation (Axiom III und IV); in $T$ sei ein
fester, das Einheitselement enthaltender Unterring $R$ ausgezeichnet. Die Begriffe
Modul, Ordnung, ganz sollen sich durchweg auf diesen festen Ring $R$ beziehen,
was der Kürze halber später in der Bezeichnung weggelassen wird.\footnote{Die
Theorie der ganzen Größen bleibt bestehen, wenn Nullteiler zugelassen werden, und
wenn dafür in $R$ der Teilerkettensatz vorausgesetzt wird, was nach §2 auch den
Beweis des Hilfssatzes ermöglicht. Da diese Fassung der ganzen Größen aber später
nicht benutzt wird, soll nur in Fußnoten darauf hingewiesen werden. Alle in Rede
stehenden Definitionen bleiben bei Existenz von Nullteilern erhalten; die meisten
benötigen — wie bekannt und leicht ersichtlich — noch geringere Voraussetzungen.
Nicht erhalten bleibt die Dedekindsche Fassung: „Eine Zahl ist ganz, wenn sie eine
Hülle besitzt.“ Denn beim Auftreten von Nullteilern braucht der Quotient zweier
endlicher Moduln nicht endlich zu bleiben; deshalb ist hier auch die Ordnung direkt
und nicht als Modulquotient definiert.} Die Elemente aus $R$ seien mit lateinischen,
die aus $T$ allgemein mit griechischen Buchstaben bezeichnet.

1. Ein System $M$ von Elementen aus $T$ heißt wie üblich $R$-Modul — kurz
Modul —, wenn es neben zwei Elementen $\alpha$ und $\beta$ auch die Differenz,
neben $\alpha$ auch $r\alpha$ enthält, unter $r$ ein beliebiges Element aus $R$
verstanden. $M$ heißt durch $N$ teilbar,
\[
        M\equiv 0 \pmod N,
\]
und $M$ ein Vielfaches, $N$ ein Teiler, wenn jedes Element von $M$ auch in $N$
enthalten ist. $R$-Moduln, deren Elemente alle zu $R$ gehören, heißen Ideale in $R$.

Offenbar ist der Durchschnitt — das kleinste gemeinsame Vielfache $[\ldots,M_i,\ldots]$
— beliebig vieler Moduln aus $T$ wieder ein Modul; ebenso ist $T$ Modul. Somit
existiert eindeutig der aus einem beliebigen System $\Sigma$ von Elementen aus $T$
abgeleitete Modul als Durchschnitt aller Moduln, die $\Sigma$ umfassen. Die Summe
— der größte gemeinsame Teiler $(\ldots,M_i,\ldots)$ — eines beliebigen Systems von
Moduln ist der aus der Vereinigungsmenge abgeleitete Modul. Das Produkt $MN$
zweier Moduln ist definiert als der aus den Elementen $\mu\nu$ abgeleitete Modul,
wo $\mu$ alle Elemente aus $M$, $\nu$ alle aus $N$ durchläuft. Das Produkt genügt
also dem assoziativen und kommutativen Gesetz, da dies für die Ringelemente gilt.
Aus dem distributiven Gesetz in $T$ folgt weiter für Moduln das distributive Gesetz
\[
        (\ldots,M_i,\ldots)N=(\ldots,M_iN,\ldots),
\]
da es sich nach eben diesem Gesetz in $T$ rechts und links um den aus allen
Elementen $\mu_i\nu$ abgeleiteten Modul handelt.

Da $R$ selbst Modul ist, läßt sich die Bedingung, daß $R$ Multiplikatorenbereich
ist, ausdrücken durch $RM\equiv0\pmod M$. Da $R$ nach Voraussetzung das
Einheitselement enthält, gilt auch $M\equiv0\pmod{RM}$ und damit $M=RM$.

Ein Modul $M$ heißt endlich, wenn es mindestens ein aus endlich vielen Elementen
bestehendes System $\Sigma$ gibt, aus dem $M$ abgeleitet ist; die Elemente
$\mu_1,\ldots,\mu_n$ von $\Sigma$ heißen eine Modulbasis von
\[
        M=(\mu_1,\ldots,\mu_n).
\]
Summe und Produkt von zwei — und damit von endlich vielen — endlichen Moduln
sind wieder endlich, da sie aus der Vereinigungsmenge bzw. den Produkten der
Basiselemente abgeleitet sind. Letzteres folgt aus der Darstellung
$r_1\mu_1+\cdots+r_n\mu_n$ aller Elemente $\mu$ aus $M$ durch die Basis und aus
dem kommutativen Gesetz der Multiplikation in $T$.

2. Ein in $T$ gelegener Erweiterungsring von $R$ — der also alle Elemente von
$R$ umfaßt — heißt eine $R$-Ordnung, kurz Ordnung. Der Durchschnitt beliebig
vieler Ordnungen aus $T$ ist wieder eine Ordnung, ebenso ist $T$ Ordnung; es
existiert also eindeutig die aus einem beliebigen System $\Sigma$ von Elementen
aus $T$ abgeleitete Ordnung. Summe und Produkt von Ordnungen lassen sich
entsprechend wie für Moduln definieren; insbesondere wird $O^2=O$ für jede
Ordnung.

Jede Ordnung ist zugleich $R$-Modul; eine Ordnung heißt endlich, wenn sie
endlicher Modul ist. Da Summe und Produkt durch Bildung von Modulsumme
bzw. Produkt entstehen, werden nach 1. Summe und Produkt endlicher Ordnungen
wieder endliche Ordnungen; weiter gilt das transitive Gesetz: Ist $A$ endliche
$R$-Ordnung, $B$ endliche $A$-Ordnung, so ist $B$ auch endliche $R$-Ordnung.
Denn $B$ wird zugleich $R$-Ordnung, also $R$-Modul. Bildet aber
$\alpha_1,\ldots,\alpha_m$ eine Modulbasis von $A$ in bezug auf $R$, und
$\beta_1,\ldots,\beta_n$ eine solche von $B$ in bezug auf $A$, so bilden
ersichtlich die Produkte $\alpha_i\beta_j$ eine Modulbasis von $B$ in bezug auf $R$.

3. Ein Element $\alpha$ aus $T$ heißt ganz in bezug auf $R$, kurz ganz, wenn die
aus $\alpha$ abgeleitete Ordnung $R_\alpha$ endlich ist — anders ausgedrückt, wenn
die Teilerkette der Moduln
\[
        U_n=(\alpha^0,\alpha,\ldots,\alpha^{n-1})
\]
im Endlichen abbricht, $U_n=U_{n+1}=\cdots$.

Die Definition läßt auch die in 4. zu benutzende zweite Fassung zu: Ein Element
$\alpha$ aus $T$ heißt ganz, wenn es mindestens einen vom Nullmodul verschiedenen
Hauptmodul $C$ — d.\,h. $C$ ist aus einem Element $\gamma\ne0$ abgeleitet — gibt,
so daß das Produkt der Ordnung $R_\alpha$ mit $C$ ein endlicher Modul ist; anders
ausgedrückt, daß die Kette der Moduln $CU_n$ im Endlichen abbricht.\footnote{Sind
in $R$ Nullteiler zugelassen, so muß die Existenz eines regulären Hauptmoduls
gefordert werden; d.\,h. $\gamma$ muß als Nicht-Nullteiler, als reguläres Element
vorausgesetzt werden, was für Ringe ohne Nullteiler auf $\gamma\ne0$ zurückkommt.}

Schließlich ist die Definition auch identisch mit der üblichen Fassung: Ein Element
$\alpha$ aus $T$ heißt ganz, wenn es mindestens einer Gleichung genügt:
\[
        \alpha^n+r_1\alpha^{n-1}+\cdots+r_n=0,
\]
wo die $r$ Elemente aus $R$ sind.

Die verschiedenen Definitionen sind tatsächlich identisch. Denn da $R_\alpha$
übereinstimmt mit dem aus allen Potenzen von $\alpha$ abgeleiteten Modul, läßt
sich für endliches $R_\alpha$ immer eine Modulbasis aus diesen Potenzen wählen.
Ist eine solche etwa gegeben durch $e=\alpha^0,\alpha,\ldots,\alpha^{n-1}$, so
bedeutet das, daß die Kette der $U_n$ mit $U_n$ abbricht, und umgekehrt ergibt
sich aus $U_n=U_{n+1}=\cdots$ diese Basis. Das Abbrechen der Kette der $U_n$
ist aber auch identisch mit dem Bestehen der Gleichung
\[
        \alpha^n+r_1\alpha^{n-1}+\cdots+r_n=0.
\]
Mit $R_\alpha$ ist auch $CR_\alpha$ endlich, also bricht die Kette der $CU_n$ ab;
anderseits folgt aus der Endlichkeit von $CR_\alpha$, also dem Abbrechen der
$CU_n$, auch das Abbrechen der $U_n$ und damit die Endlichkeit von $R_\alpha$.
Denn da $C$ als vom Nullmodul verschiedener Hauptmodul vorausgesetzt ist,
ergibt $CU_n=CU_{n+1}$ auch $U_n=U_{n+1}$.

Die Ringeigenschaft der ganzen Größen und das transitive Gesetz der ganzen
Abhängigkeit beruhen auf dem

\medskip
\noindent\textbf{Hilfssatz.} Ist $O$ eine endliche Ordnung aus $T$, $\alpha$ ein
beliebiges Element aus $O$, so ist die aus $\alpha$ abgeleitete Ordnung $R_\alpha$
endlich — m.\,a.\,W.: alle Elemente einer endlichen Ordnung sind ganz.

Der Beweis ergibt sich nach den üblichen Schlüssen. Ist $\omega_1,\ldots,\omega_n$
eine Modulbasis von $O$, so gehört wegen der Ringeigenschaft auch
$\alpha\omega_i$ zu $O$. Also kommt
\[
        \alpha\omega_i=r_{i1}\omega_1+\cdots+r_{in}\omega_n,
\]
und daraus
\[
        \det(r_{ij}-e_{ij}\alpha)=0,
\]
mit $e_{ij}=0$ für $i\ne j$, $e_{ii}=e$, womit $\alpha$ als ganz nachgewiesen ist.
Für das Verschwinden der Determinante wird benutzt, daß $T$ und damit $O$ ohne
Nullteiler vorausgesetzt ist.\footnote{Ist in $R$ der Teilerkettensatz vorausgesetzt,
sind aber Nullteiler zugelassen, so wird der Hilfssatz eine unmittelbare Folge des
Modulsatzes in §2, demzufolge sich der Kettensatz auf die Moduln aus $O$
überträgt. Auch dieser Beweis stammt für den Zahlkörper von Dedekind
(4. Aufl., §173, II) und stellt die erste Anwendung von Kettensätzen in der
Literatur dar. In den unter 4. zu gebenden Folgerungen benutzt Dedekind statt
dessen noch die speziellere Tatsache, daß die Anzahl der Restklassen nach einem
Ideal im algebraischen Zahlkörper endlich ist.}

Das System $G$ aller ganzen Größen aus $T$ bildet eine Ordnung. $G$ umfaßt
$R$; denn die Elemente aus $R$ sind ganz, da $R$ einen endlichen $R$-Modul mit
dem Einheitselement als Basis bildet. $G$ hat aber auch Ringeigenschaft; denn die
aus zwei beliebigen ganzen Größen $\alpha,\beta$ abgeleitete Ordnung $R_{\alpha,\beta}$
wird gleich dem Produkt $R_\alpha R_\beta$, also endlich. Somit sind nach dem
Hilfssatz $\alpha+\beta$ und $\alpha\beta$ als Elemente einer endlichen Ordnung ganz.

\medskip
\noindent\textbf{Transitives Gesetz der ganzen Abhängigkeit.} Ist $O$ eine aus
ganzen Größen bestehende Ordnung, so ist jedes in bezug auf $O$ ganze Element
aus $T$ auch ganz in bezug auf $R$. Nach Definition genügt $\alpha$ einer Gleichung
\[
        \alpha^m+\omega_1\alpha^{m-1}+\cdots+\omega_m=0,
\]
ist also auch ganz in bezug auf die aus $\omega_1,\ldots,\omega_m$ abgeleitete
Ordnung $O'=R_{\omega_1,\ldots,\omega_m}$, die als Produkt
$R_{\omega_1}\cdots R_{\omega_m}$ endlich wird. Nach dem unter 2. gegebenen
transitiven Gesetz wird somit die aus $\alpha$ abgeleitete endliche $O'$-Ordnung
$O'_\alpha$ auch eine endliche $R$-Ordnung, und $\alpha$ wird nach dem Hilfssatz
ganz als Element einer endlichen Ordnung.

4. Der Ring $R$ heißt ganz-abgeschlossen in $T$, wenn jedes in bezug auf $R$ ganze
Element aus $T$ zu $R$ gehört. Nach der zweiten Fassung der Definition der ganzen
Größen in 3. gehört also ein Element $\alpha$ aus $T$ dann und nur dann zu $R$,
wenn es mindestens einen vom Nullmodul verschiedenen Hauptmodul $C$ gibt, so
daß $CR_\alpha$ einen endlichen Modul bildet.\footnote{Sind in $R$ Nullteiler
zugelassen, so muß $C$ wieder als regulärer Hauptmodul vorausgesetzt werden;
ebenso muß das Element $c$ in Folgerung I als regulär vorausgesetzt werden.}

Aus dem Begriff des ganz-abgeschlossenen Ringes hat Dedekind (3. Aufl., §172)
zwei wichtige Folgerungen gezogen, die es ermöglichen, diesen Begriff für die
Idealtheorie zu verwerten:

\medskip
\noindent\textbf{Dedekindsche Folgerung I.} $R$ sei ganz-abgeschlossen in $T$,
und außerdem sei als weitere Voraussetzung in $R$ der Teilerkettensatz für Ideale
(Axiom I) erfüllt. Ist $\beta$ ein nichtganzes Element aus $T$, $c\ne0$ ein Element
aus $R$, so gibt es einen Exponenten $\rho>0$ derart, daß in der Reihe der Elemente
\[
        c,\;c\beta,\ldots,c\beta^\nu,\ldots
\]
die Elemente $c,\ldots,c\beta^\rho$ ganz, alle übrigen nichtganz sind.

Wären alle Elemente $c\beta^\nu$ ganz, so gehörten sie wegen der ganzen
Abgeschlossenheit von $R$ zu $R$. Damit aber ginge die Kette der Moduln
$C,CB_1,\ldots,CB_\nu,\ldots$ — wo $C$ den aus $c$, $B_\nu$ den aus
$1,\beta,\ldots,\beta^\nu$ abgeleiteten Modul bezeichnet — über in eine Kette von
Idealen $C_\nu=(c,c\beta,\ldots,c\beta^\nu)$, die nach Voraussetzung im Endlichen
abbricht. Damit aber wäre gegen die Voraussetzung $R_\beta$ endlich und $\beta$
ganz; es gibt somit nichtganze unter den Elementen $c\beta^\nu$. Weiter sind mit
irgendeinem nichtganzen Element auch alle folgenden nichtganz; denn ist
$c\beta^r$ ganz, also in $R$, so genügt $c\beta^\rho$ für $\rho<r$ der Gleichung
\[
        (c\beta^\rho)^r-c^{\,r-\rho}(c\beta^r)^\rho=0
\]
und ist also auch ganz. Ist somit $c\beta^{\rho+1}$ das erste nichtganze Element,
so wird — da $c$ ganz — $\rho>0$ und ist der gesuchte Exponent.

Spezialisiert man den Erweiterungsring $T$ zu dem Quotientenkörper von $R$ —
d.\,h. zu dem durch Adjunktion aller Elementenpaare entstehenden Körper\footnote{
In bekannter Steinitzscher Fassung, J. f. M. 137. Sind in $R$ Nullteiler zugelassen,
so muß an Stelle des Quotientenkörpers der Quotientenring treten, durch Adjunktion
aller Quotientenpaare, bei denen der Nenner ein reguläres Element aus $R$ ist.
Hier ist die Dedekindsche Folgerung II nicht mehr erfüllt; denn $n$ wird im
allgemeinen kein reguläres Element sein.} —, so folgt daraus die

\medskip
\noindent\textbf{Dedekindsche Folgerung II.} $R$ sei ganz abgeschlossen in seinem
Quotientenkörper $K$, und in $R$ sei der Teilerkettensatz für Ideale erfüllt. Ist
$\beta$ ein nichtganzes Element aus $K$, so läßt $\beta$ eine Darstellung als
Quotient von Elementen aus $R$ zu:
\[
        \beta=m/n
\]
derart, daß auch $m^2/n$ nichtganz ist.

Setzt man $\beta=b/c$ mit $c\ne0$, so werden in der Reihe
$c,c\beta=b,c\beta^2,\ldots$ die beiden ersten Elemente sicher ganz, also $\rho>1$,
wo $\rho$ den Exponenten aus Folgerung I bedeutet. Setzt man
\[
        m=c\beta^\rho,\qquad n=c\beta^{\rho-1},
\]
so werden $m/n=\beta$ und $m^2/n=c\beta^{\rho+1}$ nichtganz.

Das transitive Gesetz der ganzen Abgeschlossenheit besteht in der Form: Ist $R$
ein beliebiger Ring aus $T$, bezeichnet $G$ das System aller in bezug auf $R$
ganzen Größen aus $T$, so besteht $G$ auch aus allen in bezug auf $G$ ganzen
Größen aus $T$. Denn jede in bezug auf $G$ ganze Größe ist auch ganz in bezug
auf $R$, nach dem transitiven Gesetz in 3., gehört also zu $G$.

\subsection*{§2. Kettensätze in endlichen Modulbereichen}

Der Modulsatz, demzufolge sich die Kettensätze übertragen, tritt in der hier zu
gebenden Anwendung nur für Moduln eines Erweiterungsringes auf. Da der Beweis
aber im Fall eines allgemeinen Modulbereichs derselbe bleibt, soll ein solcher
zugrunde gelegt werden.

Ein System $M$ von Elementen $\alpha,\beta,\ldots$ heißt Modulbereich in bezug
auf einen Ring $R$, wenn in $M$ zwei Operationen gegeben sind, die eindeutig zu
Elementen aus $M$ führen: eine additive Verknüpfung der Elemente und eine
Multiplikation der Elemente mit den Elementen aus $R$; wenn $M$ gegenüber der
Addition eine Abelsche Gruppe bildet, während für die Multiplikation das
assoziative Gesetz erfüllt ist, und wenn das distributive Gesetz in beiden Fassungen
gilt.\footnote{Vgl. \emph{Idealtheorie}, §9.}

Für einen Modulbereich bleiben offenbar alle unter §1, 1. gegebenen
Moduldefinitionen erhalten, die sich nicht auf die Multiplikation beziehen.
Insbesondere heißt $M$ ein endlicher Modulbereich, wenn es endlich viele Elemente
$\xi_1,\ldots,\xi_k$ aus $M$ gibt derart, daß $M$ gleich dem aus
$\xi_1,\ldots,\xi_k$ abgeleiteten Modul wird, also gleich dem System aller
Linearformen $r_1\xi_1+\cdots+r_k\xi_k$, wenn in $R$ ein Einheitselement
existiert, während sonst noch Zusatzglieder $n_i\xi_i$ hinzukommen.

\medskip
\noindent\textbf{Modulsatz.} Ist $M$ ein endlicher Modulbereich in bezug auf
einen kommutativen Ring $R$ mit Einheitselement, und gilt in $R$ der Teiler- bzw.
Vielfachenkettensatz der Ideale, so gilt in $M$ der Teiler- bzw.
Vielfachenkettensatz der Moduln.

Der Beweis beruht darauf, daß jedem Modul aus $M$ eindeutig ein System von
endlich vielen Idealen aus $R$ zugeordnet wird derart, daß einem echten Teiler
bzw. Vielfachen des Moduls jeweils mindestens ein echtes Teiler- bzw.
Vielfachenideal entspricht.

Ein Element aus $M$ heißt von der Länge $i$, wenn es sich auf mindestens eine
Art als Linearform in $\xi_1,\ldots,\xi_i$ darstellen läßt, während es keine
Darstellung als Linearform in $\xi_1,\ldots,\xi_{i-1}$ zuläßt. Einem beliebigen
Modul $W$ aus $M$ seien jetzt $k$ Ideale $\mathfrak a_1,\ldots,\mathfrak a_k$
aus $R$ zugeordnet: unter $W_i$ sei der Modul aller Elemente aus $W$ verstanden
von der Länge $\le i$, und $\mathfrak a_i$ bezeichne das System aller Koeffizienten
von $\xi_i$ in $W_i$. Dann folgt vorerst:

Ist $B$ ein echter Teiler von $W$, sind $\mathfrak b_1,\ldots,\mathfrak b_k$ die
$B$ zugeordneten Ideale, so ist unter den $\mathfrak b_i$ mindestens ein echter
Teiler des entsprechenden $\mathfrak a_i$. Aus $B\equiv0\pmod W$ folgt nämlich
$B_i\equiv0\pmod {W_i}$ und damit $\mathfrak b_i\equiv0\pmod{\mathfrak a_i}$ für
$i=1,2,\ldots,k$. Nach Voraussetzung muß es in $B$ nicht in $W$ enthaltene
Elemente geben, also auch solche kleinster Länge $\ell$. Dann aber wird
$\mathfrak b_\ell$ ein echter Teiler von $\mathfrak a_\ell$; denn es wird
$\mathfrak b_\ell\equiv0\pmod{\mathfrak a_\ell}$, wenn
\[
        \beta=b_1\xi_1+\cdots+b_\ell\xi_\ell
\]
ein solches Element kleinster Länge aus $B$ ist. Aus
$\mathfrak b_\ell=\mathfrak a_\ell$ folgt nämlich die Existenz eines Elementes
\[
        \alpha=a_\ell\xi_\ell+\cdots+a_1\xi_1\equiv0\pmod W,
\]
und damit die Existenz eines Elementes $\beta-\alpha\in B$, $\beta-\alpha\notin W$
von kleinerer Länge als $\ell$.

Sei jetzt eine Teiler- bzw. Vielfachenkette
\[
        W^{(1)},W^{(2)},\ldots,W^{(\nu)},\ldots
\]
vorgelegt; sei in $R$ der Teiler- bzw. Vielfachenkettensatz erfüllt. Bezeichnen
$\mathfrak a_{1\nu},\ldots,\mathfrak a_{k\nu}$ die $W^{(\nu)}$ zugeordneten
Ideale, so bilden die Reihen
\[
        \mathfrak a_{i1},\mathfrak a_{i2},\ldots
\]
Teiler- bzw. Vielfachenketten für $i=1,\ldots,k$; es gibt also nach Voraussetzung
einen Index $\nu_0$ derart, daß das Ideal $\mathfrak a_{i\nu_0}$ gleich allen
folgenden wird für jedes $i$. Dann aber wird der Modul $W^{(\nu_0)}$ nach dem
oben Bewiesenen gleich allen folgenden, womit die Kettensätze übertragen sind.

Unmittelbar ergibt sich jetzt die

\medskip
\noindent\textbf{Folgerung des Modulsatzes.} Gilt in $R$ der
Vielfachenkettensatz modulo jedem vom Nullideal verschiedenen Ideal, so gilt in
$M$ der Vielfachenkettensatz modulo jedem Modul $C$, dessen zugeordnete Ideale
$\mathfrak c_1,\ldots,\mathfrak c_k$ sämtlich vom Nullideal verschieden sind.
Denn unter diesen Annahmen brechen die Vielfachenketten
$\mathfrak a_{i1},\ldots,\mathfrak a_{i\nu},\ldots$ im Endlichen ab.

\medskip
\noindent\textbf{Zusatzbemerkung.} Ist $R$ ein Ring ohne Einheitselement, so
überträgt sich noch der Teilerkettensatz und der Vielfachenkettensatz modulo den
vom Nullideal verschiedenen Idealen. Man betrachte nämlich an Stelle von $M$ das
System $M'$ aller Elemente der Form
\[
        a_1\xi_1+\cdots+a_k\xi_k+n_1\xi_1+\cdots+n_k\xi_k,
\]
das wieder in $M$ übergeht, wenn man die zusätzlichen ganzen Vielfachen
$n_i\xi_i$ durch Elemente von $M$ ersetzt. Diesem $M'$ lassen sich entsprechend
wie oben $2k$ Ideale zuordnen, wobei die $\mathfrak a_i$ Ideale aus $R$, die
$\mathfrak n_i$ Ideale aus den ganzen Zahlen bedeuten. Da für die
$\mathfrak n_i$ der Teilerkettensatz und der Vielfachenkettensatz modulo jedem
vom Nullideal verschiedenen Ideal erfüllt ist, bleibt für $M'$ und damit für $M$
der Beweis für den Teilerkettensatz erhalten, während man beim
Vielfachenkettensatz fordern muß, daß die $C$ bzw. $C'$ zugeordneten $2k$ Ideale
alle vom Nullideal verschieden sind.

\end{document}
